% Paper 2 	extperiodcentered{}  Appendix A 	extperiodcentered{}  Drift detection details
% Status: draft 	extperiodcentered{}  2026-05-05 	extperiodcentered{}  [Chinese comment] §3.A [Chinese comment]

\appendix

\section{Drift detection 	extperiodcentered{}  technical details}
\label{app:drift}

This appendix expands on the orthogonal anchor-based drift detection
introduced in §\ref{sec:method} and §\ref{sec:related}. The detector is
production-deployed in the \texttt{compass} hook and is not strictly
required for LongMemEval-S evaluation, but it is part of why we believe
\compass{} differs categorically from prior memory systems.

\subsection{Definition}

For a user prompt $q$ at hook time, we compute:

\begin{align}
\text{align}(q) &= \max_{a \in A_+} \cos(\bgem(q), \bgem(a)) \\
\text{deviate}(q) &= \max_{a \in A_-} \cos(\bgem(q), \bgem(a)) \\
\text{drift\_score}(q) &= \text{deviate}(q) - \text{align}(q)
\end{align}

where $A_+$ is a set of 25 \emph{positive anchors} (sentences that
exemplify aligned behavior at this user's domain) and $A_-$ is a set of
35 \emph{negative anchors} (sentences that exemplify drift exemplars).
Drift is flagged when $\text{drift\_score}(q) > \tau$ \emph{or} when any
single negative anchor's cosine exceeds $\theta$.

The thresholds $\tau, \theta$ are calibrated per-domain (see
§\ref{app:calib}) but typical values are $\tau = 0.05$ and $\theta = 0.78$.

\subsection{Anchor design}

Both anchor sets are written as \emph{task-style sentences} matching the
distribution of real prompts. We deliberately avoid abstract maxims
(``always be careful''); pilot studies showed that maxim-style anchors
yield AUC$\approx$0.51 (chance) while task-style anchors yield 0.92.

Example positive anchor (engineering domain, \texttt{anchors.json},
translated from the original Chinese for typesetting):
\begin{quote}\itshape
  ``ssh into cloud and inspect the PG \texttt{agent\_wallets} balance
  trajectory directly $\cdot$ do not rely on inference''
\end{quote}

Example contrast pair (engineering domain):
\begin{quote}\itshape
  ``Run the verify command and inspect the actual output $\cdot$ do not
  declare done without seeing OK'' \emph{(positive)}\\
  ``Claimed done without verifying $\cdot$ said OK without inspecting
  actual output'' \emph{(negative $\cdot$ same surface form)}
\end{quote}

The contrast pair makes the cosine geometry separate the right
behaviors despite high lexical overlap.

\subsection{Layered anchor inheritance (\#6 fusion)}

Production anchors are loaded from three layers:

\begin{enumerate}
\item \texttt{anchors\_platform\_base.json}: 15 positive + 25 negative
  generic platform anchors (e.g., ``don't claim done before
  verifying''). Maintained by \compass{} platform admins; affects all tenants.

\item \texttt{anchors\_<domain>.json}: domain-specific (e.g.,
  \texttt{anchors\_finance.json} for financial reasoning, \texttt{anchors\_legal.json}
  for legal review). Maintained by domain experts.

\item \texttt{anchors.json}: tenant/user custom anchors, refined from
  feedback signal.
\end{enumerate}

Total layered cost (with all three): 65 positive + 85 negative anchors,
about 32\,KB JSON, 200\,ms cold load (one-time on daemon start).

\subsection{Calibration}
\label{app:calib}

Thresholds $\tau, \theta$ are calibrated per-domain on a held-out
50-prompt validation set:

\begin{itemize}
\item Compute drift\_score on each prompt.
\item Sweep $\tau \in [-0.1, 0.2]$ and $\theta \in [0.5, 0.95]$.
\item Pick $(\tau, \theta)$ maximizing $F_1$ at fixed false-positive
  rate $\le 5\%$.
\end{itemize}

Per-domain thresholds (engineering domain shown):
$\tau = 0.045$, $\theta = 0.781$, AUC = 0.92, $F_1 = 0.84$,
fpr = 0.04, latency p95 = 47\,ms.

\subsection{Self-audit at session end}

In addition to per-prompt drift detection, \compass{} v0.9 adds
\emph{post-session self-audit}: after each Claude Code session,
\texttt{session\_writer.py} calls DeepSeek V3.2 to summarize the
session and assign one of \texttt{green/yellow/red} drift labels
plus 0--3 \texttt{drift\_signals} (specific evidence sentences).

Self-audit is not the same as anchor-based detection:

\begin{itemize}
\item Anchor-based: real-time at hook 	extperiodcentered{}  per-prompt 	extperiodcentered{}  cosine over
  bge-m3.
\item Self-audit: post-hoc 	extperiodcentered{}  per-session 	extperiodcentered{}  LLM judgment.
\end{itemize}

The two are complementary: anchor catches drift as it happens;
self-audit catches drift the user didn't flag. Together they form
the data feeding the \texttt{drift\_history} timeline tool.

\subsection{Performance targets in production}

For production deployment (compass.nautilus.social), drift detection
must satisfy:

\begin{itemize}
\item Latency: hook overhead $\le$ 50\,ms p95 (else affects
  perceived AI responsiveness).
\item Memory: anchor matrix fits in 200\,MB (1024-d float16 $\times$
  150 anchors $\times$ 2 sets $\sim$300\,KB raw, expand for batched
  inference).
\item False positive rate: $\le 5\%$ at chosen threshold (else users
  ignore alerts).
\item Catch rate on synthetic drift injection: $\ge 80\%$.
\end{itemize}

All four are met by the bge-m3 daemon configuration described in
\S\ref{sec:method}.

\subsection{Limitations of drift detection}

Anchor-based drift is \emph{black-box}: we don't know if the LLM is
internally drifting, only that its outputs match patterns we've labeled
as drift. White-box methods like Persona Vectors
\citep{anthropic2025persona} require model weight access we don't have.

Drift detection is also \emph{closed-set}: it only catches drift that
matches some pre-existing negative anchor. Novel drift modes require
new anchors. We mitigate by (a) deploying a feedback signal pathway
(\texttt{compass.feedback\_log} MCP tool) and (b) periodic anchor
retraining on accumulated signals (\texttt{python feedback.py
retrain}).

Finally, drift detection is currently English+Chinese-centric (bge-m3
training distribution). Other languages may show degraded AUC; we
do not have rigorous data here yet.
